
%(BEGIN_QUESTION)
% Copyright 2011, Tony R. Kuphaldt, released under the Creative Commons Attribution License (v 1.0)
% This means you may do almost anything with this work of mine, so long as you give me proper credit

Det er vanlig å plassere en bruddskive (rupture disk) mellom prosesstank og trykksikkerhetsventil (PSV) når de brukes sammen.  Forklar hensikten med å plassere bruddskiven der, og hvorfor det kan være en svært dårlig idé å plassere PSV mellom bruddskiven og tanken.

\underbar{file i00196}
%(END_QUESTION)





%(BEGIN_ANSWER)

Bruddskiver plasseres noen ganger mellom prosesstank og PSV for å hindre diffuse utslipp (siden alle ventiler lekker noe i normal drift, mens en intakt bruddskive ikke skal lekke).  Samme lekkasjeproblem forklarer hvorfor man aldri bør plassere en bruddskive på utløpet av en trykksikkerhetsventil.

\vskip 10pt

For å illustrere problemet: tenk en PSV som åpner ved 800 PSI, og en bruddskive med samme sprengtrykk plassert på utløpet av PSV-en.  Anta så at PSV-en lekker litt over tid, slik at røret mellom PSV og bruddskive gradvis bygges opp til 500 PSI damp.  Hvor mye prosesstrykk kreves da for å løfte PSV-en?  Husk at en selvvirkende PSV reagerer på \textit{trykkforskjellen} mellom innløp og utløp!

%(END_ANSWER)





%(BEGIN_NOTES)


%INDEX% Safety, overpressure protection: rupture disk
%INDEX% Safety, overpressure protection: safety valve

%(END_NOTES)

